% -----------------------------------------------
% DCC831-TP1: Geração Simbólica de Melodias no Estilo Coral de Bach com Cadeias
%             de Markov de Ordem n Erick Soares de Souza, UFMG 2026/2
% -----------------------------------------------

\documentclass{article}
\usepackage[T1]{fontenc}
\usepackage[utf8]{inputenc}
\usepackage{ismir}
\usepackage{amsmath,cite,url}
\usepackage{graphicx}
\usepackage{color}
\usepackage{booktabs}

\title{Geração Simbólica de Melodias no Estilo Coral de Bach com Cadeias de Markov de Ordem $n$}

\oneauthor
  {Erick Soares de Souza} {Universidade Federal de Minas Gerais \\
  \texttt{ericksoares@dcc.ufmg.br}}

\def\authorname{E. S. Souza}

% ismir.sty define \abstract com o texto fixo "ABSTRACT" em inglês;
% sobrescrevemos aqui pra manter o resumo todo em português.
\renewcommand{\abstract}{\begin{center}\bf RESUMO\end{center}}

\sloppy

\begin{document}

\maketitle

% -----------------------------------------------
\begin{abstract}
  Este trabalho apresenta um sistema de geração simbólica de melodias baseado em
  Cadeias de Markov de ordem $n$, treinado sobre a voz do baixo dos 371 corais a
  quatro vozes (SATB) de J.S. Bach distribuídos com a biblioteca
  \texttt{music21}, adotados como restrição estilística. Entre as quatro vozes,
  o baixo é escolhido por ser a mais ativa, com maior salto melódico médio e
  maior densidade rítmica, e todas as peças são normalizadas para Dó maior antes
  do treino, de modo a concentrar a aprendizagem nos padrões melódicos relativos
  ao estilo. Os tokens representam eventos de pitch e duração no formato
  \texttt{Pitch:dur}, além de pausas (\texttt{R:dur}) e marcadores de compasso
  (\texttt{BAR}). Dois hiperparâmetros principais são explorados: a ordem $n \in
    \{1, 2, 3\}$, que define a memória do modelo, e a temperatura $T \in
    \{0{,}5;\; 1{,}0;\; 1{,}5\}$, que controla a diversidade da amostragem. Para
  avaliar objetivamente a fidelidade estilística, as 9 combinações resultantes
  são comparadas estatisticamente, com 5 sementes cada, ao corpus real por meio
  de entropia de pitch, proporção de movimento por grau conjunto, intervalo
  médio e divergência de Kullback-Leibler (KL) da distribuição de intervalos. Os
  resultados indicam que a ordem 2 com $T=1{,}0$ é a configuração mais fiel ao
  estilo de Bach, com a menor divergência KL entre todas ($0{,}022$),
  equilibrando fidelidade estatística e diversidade. O código está disponível em
  \url{https://github.com/ericksoares13/DCC831-TP1}.
\end{abstract}

% -----------------------------------------------
\section{Introdução}\label{sec:intro}

A geração algorítmica de música simbólica, isto é, a geração de sequências de
eventos musicais discretos como notas e pausas, é um problema clássico de
inteligência artificial aplicada à música, discutido por Ames
(1989)~\cite{Ames1989:markov}. Entre os modelos probabilísticos propostos para
essa tarefa, as Cadeias de Markov destacam-se pela interpretabilidade, pela
facilidade de treinamento e pela capacidade de produzir resultados musicalmente
inteligíveis mesmo em ordens baixas.

Neste trabalho, adota-se o estilo coral de Bach como restrição estilística,
explorando um corpus de grande porte e estilisticamente coeso, o conjunto de 371
corais a quatro vozes distribuído com a biblioteca \texttt{music21}. A voz do
baixo, detalhada na Seção~\ref{sec:method}, é utilizada como linha de
treinamento, com normalização tonal aplicada para reduzir o viés de tonalidade
original de cada peça.

O sistema proposto é inspirado no trabalho de Shapiro \& Huber
(2021)~\cite{Shapiro2021:markov}, estendido para ordem $n > 1$, amostragem com
temperatura e um corpus de maior escala, 371 corais, contra um único arquivo por
instrumento em Shapiro \& Huber. Diferentemente da avaliação qualitativa desses
autores, este trabalho também conduz uma avaliação quantitativa sistemática
frente ao corpus real (Seção~\ref{sec:results}).

% -----------------------------------------------
\section{Trabalhos Relacionados}\label{sec:related}

Ames (1989)~\cite{Ames1989:markov} fornece o panorama fundacional do uso de
Cadeias de Markov como modelo composicional. O autor classifica os modelos por
ordem e discute o compromisso entre memorização do corpus e generalização
criativa, argumento que guia diretamente a escolha de ordem adotada. O artigo
demonstra que já na ordem 1 emergem estruturas musicalmente reconhecíveis,
enquanto ordens elevadas tendem à reprodução literal de trechos do corpus.

Shapiro \& Huber (2021)~\cite{Shapiro2021:markov} são a referência técnica
direta deste trabalho. Os autores implementam um gerador de ordem 1 treinado
sobre trechos de MusicXML de um único instrumento por corpus. A abordagem aqui
apresentada expande esse método para ordem $n$, temperatura controlável e um
corpus de maior escala (371 corais de Bach), mantendo a mesma estrutura de
Parser/Generator e os mesmos nomes de funções.

Pachet (2003)~\cite{Pachet2003:continuator} apresenta o \emph{Continuator},
sistema interativo de improvisação que aprende o estilo do músico em tempo real
via Cadeias de Markov de ordem variável, navegando o espaço de estados sem fixar
a ordem \emph{a priori} e evitando o compromisso ordem-generalização.
Diferentemente, a abordagem adotada usa ordem fixa, mais simples de inspecionar
e justificar em um contexto didático.

Conklin \& Witten (1995)~\cite{Conklin1995:multiple} propõem os \emph{Multiple
  Viewpoint Systems}, nos quais múltiplas cadeias de Markov operam sobre
representações diferentes da música (pitch, intervalo, contorno, duração), e
suas probabilidades são combinadas. Essa abordagem mitiga a esparsidade de
estados de ordem elevada, problema observado experimentalmente mesmo num
corpus estilisticamente homogêneo como o dos corais (Seção~\ref{sec:results}).
Uma extensão natural do método proposto seria adotar múltiplos pontos de
vista.

Pearce \& Wiggins (2004)~\cite{Pearce2004:methods} estendem o trabalho de
Conklin \& Witten com o modelo IDyOM, que seleciona a ordem adaptivamente com
base na entropia local da sequência. Esse mecanismo seria especialmente valioso
no contexto deste estudo, já que a esparsidade cresce rapidamente com a ordem
mesmo dentro de um corpus coeso, sugerindo que uma ordem adaptativa equilibraria
melhor fidelidade e generalização música a música.

% -----------------------------------------------
\section{Método}\label{sec:method}

\subsection{Conjunto de Dados e Pré-processamento}

É utilizada a voz do baixo dos 371 corais a quatro vozes de J.S. Bach
distribuídos com o \texttt{music21}, acessados diretamente via
\texttt{music21.corpus}. Essa voz é escolhida, entre as quatro, por ser a mais
ativa, apresentando salto melódico médio de 3,10 semitons e 47\,\% de notas com
duração menor que um tempo, contra 1,7--2,0 semitons e 19--42\,\% nas demais
vozes, medido sobre o dataset inteiro. Um pequeno número de corais não tem a voz
anotada por nome e, nesses casos, como todo coral a quatro vozes lista
soprano/alto/tenor/baixo nessa ordem, a última voz é usada como alternativa.
Cada peça é convertida com \texttt{music21}, de modo que as notas são extraídas
com seus offsets e durações expressas em \emph{quarter lengths}. Aplica-se
normalização tonal (\texttt{normalize\_key}), em que a tonalidade é detectada
por \texttt{music21.analysis.discrete} e todas as notas são transpostas para Dó
maior/lá menor antes da tokenização. Um cache em disco evita reprocessamento, já
que processar o dataset inteiro com o \texttt{music21} leva cerca de 75
segundos.

\subsection{Representação por Tokens}

Cada evento musical é mapeado para um token de texto: nota (\texttt{G4:1.0},
pitch em notação científica e duração em \emph{quarter lengths}), pausa
(\texttt{R:0.5}, silêncio com a duração indicada) e compasso (\texttt{BAR},
marcador de barra de compasso), guiando a geração pela estrutura de frases do
estilo. Como a voz do baixo é essencialmente monofônica, múltiplas notas
simultâneas praticamente não ocorrem. Por segurança, mantém-se ainda assim a
regra de reduzir qualquer sobreposição à nota mais aguda, evitando que uma
eventual sobreposição pontual exploda o espaço de estados em vez de ser tratada
como uma nota só. A sequência de tokens de cada peça constitui, assim, uma
\emph{frase de treinamento} para o modelo.

\subsection{Cadeia de Markov de Ordem $n$}

Seguindo Shapiro \& Huber (2021)~\cite{Shapiro2021:markov}, a estrutura de dados
central é um dicionário esparso de transições. Para cada estado (janela de $n$
tokens consecutivos), armazena-se a contagem de cada sucessor observado no
corpus, e após o treinamento essas contagens são normalizadas em probabilidades.

A inicialização usa a distribuição acumulada de estados iniciais (CDF) e a
função \texttt{find\_nearest\_above}, que retorna o primeiro estado cuja
probabilidade acumulada ultrapassa um valor sorteado, idêntica à de Shapiro \&
Huber (2021)~\cite{Shapiro2021:markov}. As transições seguintes usam
\texttt{random.choices}, função do Python que sorteia um elemento
proporcionalmente a pesos dados, equivalente ao método da CDF para $T = 1$.
Quando a janela atual não tem sucessor observado, mais comum em ordens altas
(Seção~\ref{sec:results}), a geração reinicia a partir da distribuição de
estados iniciais. Shapiro \& Huber, nessa situação, encerram o programa.

\subsection{Temperatura}

Antes da amostragem, os pesos de transição são reescalados:
\begin{equation}
  w'_i = w_i^{1/T}, \quad T > 0.
  \label{eq:temperature}
\end{equation}
Para $T < 1$ o modelo favorece transições mais prováveis, ficando mais
conservador, enquanto para $T > 1$ a distribuição se uniformiza, ficando mais
criativo. Em $T = 1$ o comportamento reproduz o de Shapiro \& Huber.

\subsection{Síntese MIDI}

Ao converter a sequência de tokens gerada em um arquivo MIDI, duas correções são
aplicadas. Cada pitch é restrito à tessitura do piano (21--108), transpondo por
oitavas quando necessário, já que normalização tonal ou notas atípicas do
dataset podem gerar pitches fora desse intervalo. Além disso, a velocidade MIDI
de cada nota segue um envelope senoidal por frases de 8 notas, com decaimento
global de 10\% ao longo da peça, substituindo o volume fixo (100) de Shapiro \&
Huber por uma dinâmica que sugere fraseado musical.

% -----------------------------------------------
\section{Resultados}\label{sec:results}

Treinado sobre as 371 peças, usando a voz do baixo normalizada para Dó maior/lá
menor, o vocabulário de tokens distintos é de 180 estados. O número de chaves de
transição cresce com a ordem, sendo 180 em ordem 1, 1861 em ordem 2 e 6644 em
ordem 3, o que evidencia a esparsidade crescente mesmo num corpus
estilisticamente homogêneo. A normalização tonal reduz esse número em todas as
ordens, 1861 contra 2155 chaves em ordem 2 e 6644 contra 7873 em ordem 3 sem
normalização, confirmando que concentrar as peças em uma única tonalidade
central reduz a fragmentação de padrões equivalentes transpostos.

Para avaliar a fidelidade estilística, são geradas as 9 combinações de ordem
($n\in\{1,2,3\}$) e temperatura ($T\in\{0{,}5;1{,}0;1{,}5\}$), com 5 sementes
cada (200 tokens, 80 BPM), comparadas estatisticamente com os 371 corais reais
(mesma voz, mesma normalização). As métricas usadas incluem entropia de Shannon
do pitch ($H$), proporção de intervalos de até 2 semitons (\emph{steps}, limiar
de cantabilidade), intervalo melódico médio absoluto, proporção de tokens de
pausa/compasso, divergência KL entre a distribuição de intervalos gerada e a
real, e o desvio médio em $\sigma$ (mean-$z$) em relação à variabilidade
observada entre os próprios corais. A Tabela~\ref{tab:results} resume os
resultados, com a primeira linha trazendo os valores de referência do corpus
real (média por coral).

\begin{table}[ht]
  \centering
  \footnotesize
  \setlength{\tabcolsep}{3pt}
  \caption{Métricas vs.\ corpus real (371 corais, 9 configurações, 5 sementes cada).}
  \label{tab:results}
  \begin{tabular}{ccrrrrrr}
    \toprule
    Ord                      & $T$   & $H$    & Steps  & Ivl $\mu$ & Pausas &
    $z$                      & KL                                             \\
    \midrule
    \multicolumn{2}{c}{Bach} & 4,241 & 66,1\% & 3,249  & 21,1\%    & 0,00   &
    0,000                                                                     \\
    1                        & 0,5   & 4,152  & 70,3\% & 3,097     & 25,2\% &
    0,50                     & 0,079                                          \\
    1                        & 1,0   & 5,069  & 57,1\% & 4,138     & 23,4\% &
    2,16                     & 0,149                                          \\
    1                        & 1,5   & 5,441  & 55,1\% & 4,555     & 16,7\% &
    2,85                     & 0,208                                          \\
    2                        & 0,5   & 4,575  & 70,3\% & 3,132     & 28,8\% &
    0,90                     & 0,041                                          \\
    2                        & 1,0   & 4,820  & 67,3\% & 3,129     & 22,2\% &
    1,22                     & 0,022                                          \\
    2                        & 1,5   & 5,110  & 62,7\% & 3,396     & 20,4\% &
    1,63                     & 0,025                                          \\
    3                        & 0,5   & 4,793  & 64,3\% & 3,546     & 25,9\% &
    1,31                     & 0,023                                          \\
    3                        & 1,0   & 4,978  & 62,7\% & 3,464     & 22,8\% &
    1,48                     & 0,034                                          \\
    3                        & 1,5   & 5,013  & 62,9\% & 3,598     & 23,4\% &
    1,65                     & 0,029                                          \\
    \bottomrule
  \end{tabular}
\end{table}

Em termos de fidelidade, a ordem 2 com $T=1{,}0$ é a configuração mais próxima
do corpus real pela distribuição de intervalos, com a menor divergência KL entre
todas ($0{,}022$). Seu intervalo médio, 3,13 contra 3,25 semitons reais, e sua
proporção de steps, 67,3\,\% contra 66,1\,\% reais, ficam dentro de 1$\sigma$ do
corpus, e a proporção de pausas, 22,2\,\% contra 21,1\,\% reais, é quase
idêntica, sendo a única configuração que acerta ambos simultaneamente.

Quanto à relação entre ordem e fidelidade, as ordens 2 e 3 produzem
distribuições de intervalo muito mais próximas do real que a ordem 1, cujo
$T=1{,}5$ chega a KL $=0{,}208$, nove vezes maior que o de $O2{+}T1{,}0$. Isso
ocorre porque a cadeia de ordem 1 aprende transições isoladas e, com temperatura
alta, passa a gerar saltos de 7\textordmasculine/8\textordmasculine\ que Bach
raramente usa, enquanto as ordens 2--3 preservam contexto suficiente para
restringir os saltos a regiões musicalmente plausíveis do estilo.

Já o efeito da temperatura é assimétrico entre ordens. Em ordem 1, $T=1{,}5$
causa colapso de fidelidade, com intervalo médio de 4,55 semitons e só 55\,\% de
steps, ao passo que em ordens 2--3 o efeito é suave, com o KL de ordem 2 subindo
apenas de 0,022 ($T{=}1{,}0$) para 0,025 ($T{=}1{,}5$), ainda dentro de
1$\sigma$. Isso indica que a temperatura só introduz diversidade genuína quando
há contexto suficiente, enquanto em ordem 1 ela introduz principalmente ruído.

Analisando-se a entropia, nota-se que maior fidelidade não implica menor
entropia. A configuração mais próxima globalmente por mean-$z$ ($O1{+}T0{,}5$,
mean-$z=0{,}50$) tem a menor entropia, 4,15 bits, mais perto dos 4,24 reais,
enquanto $O2{+}T1{,}0$ tem entropia maior, 4,82 bits, $2{,}0\sigma$ acima, e
ainda assim é a mais fiel pelos intervalos. A razão é que Bach usa uma
distribuição de pitch não uniforme, em que tônica, dominante e subdominante
dominam, e a ordem 2 distribui alturas de forma mais homogênea, mas preserva os
intervalos preferenciais do estilo (graus conjuntos, quartas/quintas).

Observa-se também uma saturação em ordem 3. De $T{=}1{,}0$ para $T{=}1{,}5$, o
KL \emph{cai} (0,034$\to$0,029) e a entropia varia só 0,04 bits, já que, com
contexto de 3 tokens, o espaço de transições é tão restrito que a temperatura
extra redistribui peso entre opções igualmente raras, sem acrescentar
diversidade real.

Em suma, ordem 2 com $T{=}1{,}0$ é a configuração recomendada, com KL
$=0{,}022$, ordem 1 com $T{=}0{,}5$ serve de referência conservadora, com menor
variância e entropia mais próxima da real, e ordem 1 com $T{=}1{,}5$ ilustra
divergência máxima, com KL $=0{,}208$, saltos frequentes e pouca estrutura
frásica.

% -----------------------------------------------
\section{Discussão e Trabalhos Futuros}\label{sec:discussion}

\subsection{Limitações}

A principal limitação está na natureza das métricas empregadas. Entropia,
divergência KL e mean-$z$ capturam estatísticas agregadas de pitch e intervalo,
mas não avaliam estrutura musical de mais alto nível, como frase, cadência ou
forma. Uma avaliação perceptual complementar, idealmente com ouvintes de
formação musical, ainda seria valiosa para validar essas conclusões. O autor,
sem essa formação, apoiou a análise sobretudo nas métricas objetivas por esse
motivo.

Duas outras limitações vêm da representação escolhida. O sistema gera linhas sem
harmonização, mesmo treinado sobre um repertório harmonizado, já que o baixo é
só uma das quatro vozes de cada coral, o que restringe a percepção musical do
resultado a uma única linha melódica. Já as notas simultâneas na voz do baixo
são reduzidas à nota mais aguda, embora isso seja raro nesse tipo de voz
predominantemente monofônica.

Por fim, a normalização tonal força todas as peças para Dó maior, perdendo a
identidade tonal original em favor da concentração estatística nos padrões
melódicos relativos. Uma alternativa seria normalizar apenas os intervalos,
preservando o modo.

\subsection{Trabalhos Futuros}

Uma extensão imediata é adotar a abordagem de múltiplos pontos de vista de
Conklin \& Witten (1995) e Pearce \& Wiggins
(2004)~\cite{Conklin1995:multiple,Pearce2004:methods}, combinando cadeias sobre
pitch, intervalo e duração para mitigar a esparsidade de estados de ordem
elevada. Outra direção natural, dado que o corpus já contém as quatro vozes de
cada coral, é condicionar a cadeia do baixo às demais vozes (soprano, alto,
tenor) como harmonização, análogo ao Continuator de Pachet
(2003)~\cite{Pachet2003:continuator}, mas guiado pela harmonia real de Bach, não
por interação em tempo real. Por fim, modelos de linguagem musical como Music
Transformer, Coconet ou DeepBach oferecem um patamar comparativo natural para
avaliar o ganho das abordagens neurais sobre o baseline Markoviano.

% -----------------------------------------------
\section{Declaração de Uso de IA}

Este trabalho utilizou assistência de IA generativa (Claude, Anthropic), via
\emph{Claude for Scientists}, como apoio na implementação do código, na análise
quantitativa dos resultados e na revisão do resumo. Detalhes do uso e o acesso
às conversas estão declarados no README do repositório
(\url{https://github.com/ericksoares13/DCC831-TP1}).

% -----------------------------------------------
\bibliography{ISMIRtemplate}

\end{document}
